%*******************************************************
% Abstract in English
%*******************************************************
\pdfbookmark[0]{Abstract}{Abstract}


\begin{otherlanguage}{american}
	\chapter*{Abstract}

	\begin{sloppypar}
	\noindent This thesis investigates performance optimization of an APT scanner by replacing object-based, recursively synchronous scheduling with a centralized task-based scheduling model.

	\vspace{7.56mm}

	\noindent Three scheduling algorithms were implemented and evaluated: First-In, First-Out (FIFO), Shortest-Subtree-Time-First (SSTF), and Work Stealing with SSTF prioritization. The evaluation was conducted experimentally in a controlled test environment on Debian~13 and Windows~11 with 1, 4, and 10 threads and 20 repetitions per configuration. The measured metrics were runtime, CPU and RAM utilization, waiting time, burst time, and the share of synchronously executed tasks.

	\vspace{7.56mm}

	\noindent On Windows, the new schedulers show clear runtime improvements at 4 and 10 threads compared with the reference implementation. On Linux, runtime differences are small. SSTF provides the most stable overall balance of runtime, waiting time, and memory behavior. Work Stealing does not show a consistent improvement over SSTF.

	\vspace{7.56mm}

	\noindent The results show that the architectural change is feasible and enables reproducible performance gains, especially under high-load profiles.
	\end{sloppypar}

\end{otherlanguage}
